% \iffalse meta-comment
% !TeX program  = XeLaTeX
% !TeX encoding = UTF-8
%
% Copyright (C) 2026--2026 by Innocent <InnocentFIVE@outlook.com>
%
% This work may be distributed and/or modified under the
% conditions of the LaTeX Project Public License, either
% version 1.3c of this license or (at your option) any later
% version. The latest version of this license is in:
%
%   http://www.latex-project.org/lppl.txt
%
% and version 1.3 or later is part of all distributions of
% LaTeX version 2005/12/01 or later.
%
% This work has the LPPL maintenance status `maintained'.
%
% The Current Maintainer of this work is Innocent.
%
% This work consists of the files sumcorner.dtx,
%           and the derived files sumcorner.sty,
%                                 sumcorner.pdf,
%                             and README.md.
%
%<*internal>
\iffalse
    %</internal>
    %
    %<*readme>
    README
    %</readme>
    %
    %<*internal>
\fi
\begingroup
\def\NameOfLaTeXe{LaTeX2e}
\expandafter\endgroup\ifx\NameOfLaTeXe\fmtname\else
\csname fi\endcsname
%</internal>
%
%<*install>
\input docstrip.tex
\keepsilent
\askforoverwritefalse

\preamble

Copyright (C) 2026--2026 by Innocent <InnocentFIVE@outlook.com>

This work may be distributed and/or modified under the
conditions of the LaTeX Project Public License, either
version 1.3c of this license or (at your option) any later
version. The latest version of this license is in:

http://www.latex-project.org/lppl.txt

and version 1.3 or later is part of all distributions of
LaTeX version 2005/12/01 or later.

This work has the LPPL maintenance status `maintained'.

The Current Maintainer of this work is Innocent.

This work consists of the files sumcorner.dtx,
and README.md.

\endpreamble

\generate{
    \file{\jobname.sty}        {\from{\jobname.dtx}{package}}
    %<*internal>
    \file{\jobname.ins}        {\from{\jobname.dtx}{install}}
    %</internal>
    \nopreamble\nopostamble
    \file{README.md}           {\from{\jobname.dtx}{readme}}
}

\obeyspaces

\endbatchfile
%</install>
%
%<*internal>
\fi
%</internal>
%
%<class|package>\NeedsTeXFormat{LaTeX2e}[2020/10/01]
%<*!(driver|install)>
%<!readme>\GetIdInfo $Id: sumcorner.dtx 0.3 2026-01-23 12:00:00Z Innocent <InnocentFIVE@outlook.com> $
%<package>  {sumcorner}
%<package>\ProvidesExplPackage{\ExplFileName}
%<package>{2026-01-23} {} {At the upper right corner of summation}
%<readme>人类从未知晓萨姆之角的真实样貌。
%<readme>
%<readme>> 多年以后，翻开*The TeXbook*，你仍然将会回想起我带你去见识萨姆之角的那个遥远的下午。
%<readme>
%<readme># 萨姆之角
%<readme>
%<readme>这是一个尝试，用来证明萨姆之角存在性的尝试，为此，我们大量地使用了幻影。除了在文档之中，我们很难意识到萨姆之角的存在性：这是一个非常微妙的细节问题。
%<readme>为了阐明清楚这个概念，我们不得不从1991年开始说起。著名计算机科学家高德纳在他的著作*The TeXbook*中提出萨姆之角这一概念（事实上，是习题18.44），他将其视为一个难度等级为$2$的习题。
%<readme>之后，他在同一本著作中，给出了如下论证：
%<readme>
%<readme>```tex
%<readme>\def\sumprime_#1{\setbox0=\hbox{$\scriptstyle{#1}$}
%<readme>\setbox2=\hbox{$\displaystyle{\sum}$}
%<readme>\setbox4=\hbox{${}'\mathsurround=0pt$}
%<readme>\dimen0=.5\wd0 \advance\dimen0 by-.5\wd2
%<readme>\ifdim\dimen0>0pt
%<readme>\ifdim\dimen0>\wd4 \kern\wd4 \else\kern\dimen0\fi\fi
%<readme>\mathop{{\sum}'}_{\kern-\wd4 #1}}
%<readme>```
%<readme>
%<readme>这个论证正是我们的雏形。
%</!(driver|install)>

%<*driver>
\documentclass{fdudoc}
\usepackage{sumcorner,xpatch}
\usepackage{makeidx}

\ExplSyntaxOn

\xpatchcmd \__codedoc_print_macroname:nN
{ \__codedoc_get_hyper_target:xN }
{ \__codedoc_get_hyper_target:eN }
{ } { \patchfailed }
\xpatchcmd \__codedoc_print_macroname:nN
{
    \tl_replace_all:Non \l__codedoc_tmpa_tl
    { \c_catcode_other_space_tl }
    { \fontspec_visible_space: }
    \__codedoc_macroname_prefix:o \l__codedoc_tmpa_tl
    \__codedoc_macroname_suffix:N #2
}
{ \tl_replace_all:NVn \l__codedoc_tmpa_tl
    \c_catcode_other_space_tl
    { \fontspec_visible_space: }
    \__codedoc_print_macroname_aux:on
    \l__codedoc_tmpa_tl
    { \bool_if:NT #2 { \__codedoc_typeset_TF: } } }
{ } { }
\makeatletter
\lstnewenvironment{fsyntax}[1][]{%
    \lstset{style=style@syntax, #1}\vspace{-1.5ex}}{}

\AtBeginEnvironment { fsyntax }
{
    \cs_set:Npn \lparen { \textup { ( } }
    \cs_set:Npn \rparen { \textup { ) } }
    \char_set_catcode_active:N |
    \char_set_catcode_active:N <
    \char_set_active_eq:NN | \orbar
    \char_set_active_eq:NN < \syntaxopt@aux
}
\makeatother
\ExplSyntaxOff

\newcommand{\floor}[1]{\lfloor#1\rfloor}
\setmathfont{Asana-Math.otf}

\begin{document}
\DocInput{\jobname.dtx}
\clearpage
\IndexLayout
\PrintIndex
\end{document}
%</driver>
% \fi
% \title{萨姆之角 \pkg{sumcorner} \\ \emph{At the upper right corner of summation}}
% \author{Innocent\\
% \url{InnocentFIVE@outlook.com}}
% \maketitle
% \def\d{\mathrm{dyadic}}
% \begin{abstract}
%     我们考察了萨姆之角的存在性。
% \end{abstract}
% 
% \tableofcontents
% \begin{documentation}
% \section{引论}
% 存在一些人类使用$\sum'$来表达特殊的求和，比如说二进求和之类：
% \[
%   \sumcorner\sb{n=1}^{\floor{1/\delta^N}}
%   n^p
%   \lessapprox\sb{p,N}
%   \delta^{-Np},\quad
%   \sumcorner\sb{n=1}^{A}
%   \frac{1}{n}
%   =
%   1 - O\Bigl(\frac{1}{A}\Bigr).
% \]
% 有些人类喜欢说话，因此会用$\sum^{\d}$来表示：
% \[
%   \sumcorner[\d]\sb{n=1}^{100\floor{1/\delta^N}}
%   n^p
%   \lessapprox\sb{p,N}
%   \delta^{-Np},\quad
%   \sumcorner[\d]\sb{n=1}^{A}
%   \frac{1}{n}
%   =
%   1 - O\Bigl(\frac{1}{A}\Bigr).
% \]
% 当然，我们可以更直接，使用$\sum^2$来表示，比如说：
% \[ 
%   \sumcorner[k]\sb{j=1}^{n} :=
%   \sum\sb{j\in k^{\mathbb{N}},1\leq j\leq n},\qquad
%   \sumcorner[2]\sb{j=1}^{2^m} j = \Theta(2^m),\quad
%   \sumcorner[{\sumcorner[2]\sb{j=1}^{2^m} \frac j2}]\sb{k=1}^{2^{mn}} k = \Theta(2^{n-m}).
% \]
% 当然，这会使公式稍微变长一些。
% 
% 那么，也存在一些人类会询问，如果是 |\textstyle| 情形呢？
% 
% 是的，为了防止 |\textstyle| 等情形引发歧义，我们在所有情况都需要保证 |\limits|，也就是
% $\sumcorner\sb{n=1}^{\floor{1/\delta^N}}$。
% 
% 综上所述，我们证明了萨姆之角的存在性。
% 
% \section{使用方法}
% \begin{function}{\sumcorner}
% \begin{fsyntax}
%     \sumcorner (*\oarg{*}*) _ (*\oarg{下标}*) ^ (*\oarg{上标}*)
% \end{fsyntax}
% 这样就可以了。使用 |*| 和不使用 |*| 的区别是：
% \[
% \sumcorner*^{100A+10B}\sb {j=B} f(j) ,\quad
% \sumcorner^{100A+10B}\sb {j=B} f(j) .
% \]
% \end{function}
% 
% \end{documentation}
% 
% \EnableImplementation
% \begin{implementation}
% \newgeometry{
%   left      = 2.2 in,
%   right     = 1.00 in,
%   top       = 1.25 in,
%   bottom    = 1.00 in,
%   marginpar = 2.25 in
% }
% \section{实现细节}
%    \begin{macrocode}
%<@@=sumcorner>
%<*package>
%    \end{macrocode}
%    
% 首先是检查 \LaTeX3 环境。
%    \begin{macrocode}
\RequirePackage { xtemplate, l3keys2e }
\msg_new:nnn { sumcorner } { l3-too-old }
{
    Package~ `#1'~ is~ too~ old. \\\\
    Please~ update~ an~ up-to-date~ version~ of~ the~ bundles \\
    `l3kernel'~ and~ `l3packages'~ using~ your~ TeX~ package \\
    manager~ or~ from~ CTAN.
}
\clist_map_inline:nn { xtemplate, l3keys2e }
{
    \@ifpackagelater {#1} { 2020/07/17 }
    { } { \msg_error:nnn { sumcorner } { l3-too-old } {#1} }
}
%    \end{macrocode}
% \subsection{内部变量声明}
% \begin{variable}{
% \l_@@_tmpa_box,
% \l_@@_tmpb_box,
% \l_@@_tmpc_box,
% \l_@@_tmpd_box,
% \l_@@_tmpe_box,
% \l_@@_tmpa_dim,
% \l_@@_tmpb_dim,
% \l_@@_tmpc_dim,
% \l_@@_tmpd_dim
% }
% 临时变量声明。
%    \begin{macrocode}
\box_new:N \l_@@_tmpa_box
\box_new:N \l_@@_tmpb_box
\box_new:N \l_@@_tmpc_box
\box_new:N \l_@@_tmpd_box
\box_new:N \l_@@_tmpe_box

\dim_new:N \l_@@_tmpa_dim
\dim_new:N \l_@@_tmpb_dim
\dim_new:N \l_@@_tmpc_dim
\dim_new:N \l_@@_tmpd_dim
%    \end{macrocode}
% \end{variable}
% \begin{variable}{\c_@@_math_subscript_token}
% 现在定义一个 |_| 的替代。
%    \begin{macrocode}
\tl_new:N \c_@@_math_subscript_token
\tl_set:Nn \c_@@_math_subscript_token {
  \char_generate:nn { `\_ } { 8 }
}
%    \end{macrocode}
% \end{variable}
% 
% \subsection{具体实现}
% \begin{macro}{\@@_mathpalette:nnn,\_@@_inverse_kern:N}
%    首先是一个 |\mathpalette| 的多元替代，
%    \begin{macrocode}
\cs_new:Npn \@@_mathpalette:nnn #1#2#3 {
  \tex_mathchoice:D { #1\tex_displaystyle:D { #2 } { #3 } }
  { #1\tex_textstyle:D { #2 } { #3 } }
  { #1\tex_scriptstyle:D { #2 } { #3 } }
  { #1\tex_scriptscriptstyle:D { #2 } { #3 } }
}
\cs_new:Nn \_@@_inverse_kern:N {
  \tex_kern:D #1 \tex_mathopen:D { }
}
%    \end{macrocode}
% \end{macro}
% 然后一通爆算，
% \begin{macro}{\@@_aux:}
%    \begin{macrocode}
\NewDocumentCommand { \@@_aux: }{ smmmm } {
  \hbox_set:Nn \l__sumcorner_tmpa_box { $ #3 \sum $ }
  \hbox_set:Nn \l__sumcorner_tmpb_box { $ #3 \sum \nolimits ^ { #2 } $ }
  \hbox_set:Nn \l__sumcorner_tmpc_box { $ #3 ^ { #4 } $ }
  \hbox_set:Nn \l__sumcorner_tmpd_box { $ #3 \sb { #5 } $ }
  \dim_set:Nn \l__sumcorner_tmpb_dim {
    \dim_max:nn { \box_wd:N \l__sumcorner_tmpc_box }
    { \box_wd:N \l__sumcorner_tmpd_box }
  }
  \dim_compare:nNnTF { \box_wd:N \l__sumcorner_tmpc_box } >
    { \box_wd:N \l__sumcorner_tmpa_box } {
      \IfBooleanTF { #1 } {
        \dim_set:Nn \l__sumcorner_tmpc_dim { \c_zero_dim }
      } {
        \dim_set:Nn \l__sumcorner_tmpc_dim {
          \dim_eval:n { ( \box_wd:N \l__sumcorner_tmpc_box -
            \box_wd:N \l__sumcorner_tmpa_box ) / 2 }
        }
      }
    } {
      \dim_set:Nn \l__sumcorner_tmpc_dim { \c_zero_dim }
    }
  \dim_set:Nn \l__sumcorner_tmpa_dim
  { \dim_eval:n { \box_wd:N \l__sumcorner_tmpb_box -
      \box_wd:N \l__sumcorner_tmpa_box + \l__sumcorner_tmpc_dim } }
  \dim_set:Nn \l__sumcorner_tmpd_dim
   { \dim_eval:n {
       \dim_max:nn { \box_wd:N \l__sumcorner_tmpb_box 
       - \box_wd:N \l__sumcorner_tmpa_box + \l__sumcorner_tmpc_dim  } 
       { ( \l__sumcorner_tmpb_dim - \box_wd:N \l__sumcorner_tmpa_box ) / 2  }
       - \dim_max:nn { \c_zero_dim } 
       { \l__sumcorner_tmpb_dim - \box_wd:N \l__sumcorner_tmpa_box } / 2
       - \c_zero_dim
     } } 
  \tex_mathop:D { }
  \tex_kern:D -\l__sumcorner_tmpd_dim \tex_mathopen:D { }
  \hbox_set:Nn \l__sumcorner_tmpa_box { $
      \cs_if_eq:NNTF #3 \tex_displaystyle:D { \tex_scriptstyle:D } {
        \cs_if_eq:NNTF #3 \tex_textstyle:D
        { \tex_scriptstyle:D } 
        { \tex_scriptscriptstyle:D }
      }
      #2 $ }
  \tex_mathop:D { 
    \__sumcorner_inverse_kern:N \l__sumcorner_tmpa_dim
    \sum \tex_nolimits:D \c_math_superscript_token {
      \tex_kern:D \l__sumcorner_tmpc_dim
      \box_set_ht:Nn \l__sumcorner_tmpa_box { \c_zero_dim }
      \box_use:N \l__sumcorner_tmpa_box
    } }
  \tex_limits:D
  \tl_if_blank:nTF { #4 } { } { \c_math_superscript_token { #4 } }
  \tl_if_blank:nTF { #5 } { } { \c_math_subscript_token { #5 } } 
}
%    \end{macrocode}
%  \end{macro}
%  最后简单分类讨论一下即可。
% \begin{macro}{\sumcorner}
%    \begin{macrocode}
\exp_args:NNe \NewDocumentCommand \sumcorner
{ s O{\prime} E{^\c_@@_math_subscript_token}{{}{}} t' } {
  \IfBooleanTF { #1 } {
    \IfBooleanTF { #5 } {
      \sumcorner * [ #2 ] \c_math_superscript_token
      { #3 \prime } \c_math_subscript_token { #4 }
    } {
      \@@_mathpalette:nnn
      { \@@_aux: *{ #2 } } { #3 } { #4 }
    }
  } {
    \IfBooleanTF { #5 } {
      \sumcorner [ #2 ] \c_math_superscript_token
      { #3 \prime } \c_math_subscript_token { #4 }
    } {
      \@@_mathpalette:nnn
      { \@@_aux: { #2 } } { #3 } { #4 }
    }
  }
}
%</package>
%<@@=>
%    \end{macrocode}
% \end{macro}
% \end{implementation}
% \begin{thebibliography}{1}
%   \bibitem{summation} \rule{3em}{.4pt}. 
%   ``At the upper right corner of summation,'' \emph{\rule{8em}{.4pt}}: 1--1. 2026.
% \end{thebibliography}
